% s10 The Wick tube I: torus inputs and the extensive bound
\section{The Wick tube I: torus inputs and the extensive bound}

The Wick leg controls the joint functional of the curvature and the Wick square on the tube $\Dtube$ of \cref{def:s02:tube}. Its positive leg (sources $q\ge0$) and negative leg (sources $-q$) are proved in the next section from three inputs on \emph{King's tori}: lattice reflection positivity (\cref{prop:s10:RP}), the chessboard estimate it yields (\cref{lem:s10:CB}), and an \emph{extensive bound} \Itwo\ for constant positive scalar sources in the fixed-spacing volume limit (\cref{def:s10:I2}). This section proves the first two and derives the third from the torus inputs \cref{hyp:s03:T1,hyp:s03:T2} and the sign clause \Tthree\ (\cref{prop:s10:U}); a variant with the auxiliary mass on either side of $m$ (\cref{prop:s10:Upm}) serves the discussion of the coupling range; no continuum limit is taken here.

\subsection{King's tori, the two conventions and mid-plane pairs}

\begin{definition}[King's tori]\label{def:s10:torus}
A \emph{King torus} is $T=\R^3/(\ell\Z)^3$ with $\ell\in M\N$ even; $|T|\defeq\ell^3$ is its volume, and $T_J\defeq\R^3/(2MJ\Z)^3$, $J\in\N$, is the torus with fundamental cell $\CJ$. Since $\ell\Z^3\subset\Z^3\subset\Lk$, the lattice $\Lk$ descends to $T$, with $\ell\,5^k$ sites per side, an even number, and $|T|/\eps_k^3$ sites in all. The \emph{periodic state} $\langle\cdot\rangle_{k,T}$ is King's action of Section~2 summed over the finite graph $\Lk/(\ell\Z)^3$, with periodic conditions for both fields \cite[(4.1)]{K2}, \cite[(2.21)]{K1}: Feynman-form gradient terms for every pair of parallel neighbouring bonds, the photon mass $\lmu^2$, hopping terms on every bond, the single-site weights $e^{-\eps^3[\frac12(m^2+\dm)r^2+\lambda r^4]}$. Two normalisations of the counterterm are used:
\begin{enumerate}[label=\textup{(\roman*)},nosep]
\item the \emph{free convention} of Section~2 ($\dm=\dm_\eps$, Wick constant $\sW_\eps$; the convention of every cube state of this paper), with $N^a_{k,T}(0,q)\defeq e^{-(a_k+2\sW_\eps)\sS(q)}\langle e^{\sum_x\eps^3q(x)r_x^2}\rangle_{k,T}$ for a real function $q$ on the sites of $T$, exactly as on boxes (\cref{def:s02:sources}), and
\[
\alpha^a_{k,T}(u)\defeq|T|^{-1}\ln N^a_{k,T}(0,u\mathbf 1_T)\qquad(u\in\R),
\]
the \emph{free energy density} of a constant source;
\item \emph{King's torus convention} \cite[(2.4), (3.33), (3.35)]{K1}: the counterterm $\dm_{k,T}(M)$ and the Wick constant $c_{k,T}(M)$ computed with the periodic propagators of $T$ at scalar mass $M$, together with the field-independent constants $E(M)=E_0(M)+E_1(M)$ of King's torus action. $Z^{(M)}_{k,T}\defeq\int e^{-S}$ is the source-free partition function of this model at scalar mass $M$ (bare coefficient $M^2+\dm_{k,T}(M)$; \cref{hyp:s03:T1} is a statement about these numbers), and $N^{z,\mathrm{tor}}_{k,T}(0,q)$ the normalised source functional of the torus model at mass $m$ in the reference scheme $z$, with $c_{k,T}=c_{k,T}(m)$ in place of $\sW_\eps$.
\end{enumerate}
In either convention the coefficients are the same at every site of $T$. A \emph{mid-plane pair} is $\pi=\{x_\nu=c\}\cup\{x_\nu=c+\ell/2\}$ with $\nu\in\{1,2,3\}$ and $c\in\eps_k(\Z+\frac12)$ (both planes lie between lattice planes, since $\ell5^k$ is even); $\theta$ denotes the reflection of $T$ in $\pi$, which exchanges the two slabs of $T\setminus\pi$ (``the two sides of $\pi$''), acting on the fields as in \cref{def:s07:reflection}.
\end{definition}

\subsection{Reflection positivity of King's tori}

\begin{proposition}[lattice reflection positivity of King's tori]\label{prop:s10:RP}
Let $T$ be a King torus, $k\ge0$, $\pi$ a mid-plane pair and $\theta$ the reflection in $\pi$. For every finite family $G_1,\dots,G_N$ of gauge-invariant Borel functions of the fields on one side of $\pi$ with $\langle|G_i|^2\rangle_{k,T}<\infty$ and every $\gamma\in\C^N$,
\[
\sum_{i,l=1}^N\bar\gamma_i\gamma_l\,\langle\theta(G_i)\,G_l\rangle_{k,T}\ \ge\ 0 .
\]
The periodic state is invariant under $\theta$ and under all lattice translations. This holds in both conventions of \cref{def:s10:torus}: reflection positivity does not see the value of the (site-independent) counterterm.
\end{proposition}

\begin{proof}
The proof is the mechanism of \cref{lem:s07:RPbox} with a pair of planes in place of one and without Dirichlet diagonals, in three steps.

\emph{Step 1: from the Feynman form to the Proca form.} On the torus write $A=Z+d_0\chi$ with $d_0^TZ=0$, $\chi\perp1$ (Hodge decomposition on the finite periodic graph). The Feynman form equals $\frac12\langle Z,(d_1^Td_1+\lmu^2)Z\rangle$ plus a quadratic form in $\chi$, and the Proca form $\frac12\langle A,(d_1^Td_1+\lmu^2)A\rangle$ equals the same $Z$-form plus another $\chi$-form; both $\chi$-forms are positive definite on $1^\perp$ because $\lmu>0$. For gauge-invariant $G$ the matter integral at fixed $A$ is invariant under $A\mapsto A+d_0\chi$ (rotate $\phi$ site by site), so the $\chi$-integrals factor and cancel in the normalised expectation: gauge-invariant expectations in the Feynman form equal those in the Proca form. Fubini's theorem applies by $\lmu>0$, $\lambda>0$ and a growth bound $|G|\le Ce^{a\sum_b|A(b)|}e^{C\sum_xr_x^2}$, satisfied by bounded $G$ and by the observables $\Theta(\zeta)$; so this step uses $\lmu>0$, gauge invariance, $\lambda>0$ and the growth bound, and nothing about the value of the counterterm.

\emph{Step 2: the Proca torus measure is reflection positive about $\pi$.} In the Proca form the action couples the two sides of $\pi$ only through the plaquettes and hopping terms crossing one of the two planes of $\pi$ (no gradient terms between parallel bonds across a plane; the mass term is local). Write each crossing bond variable as $a_b=\alpha_b-\alpha'_b$ with $\alpha_b,\alpha'_b$ independent centred Gaussians of variance $(2\lmu^2\eps^3)^{-1}$, attached to the two sides. Then $\theta$ becomes the swap of the two data sets $\xi=(\omega_+,\alpha)$, $\xi'=(\omega_-^\theta,\alpha')$, whose a priori measures coincide by the $\theta$-invariance of the one-sided terms, and exactly as in Step~2 of \cref{lem:s07:RPbox} every crossing plaquette term is $\frac\eps2X_p(\xi)^2+\frac\eps2X_p(\xi')^2-\eps X_p(\xi)X_p(\xi')$ and every crossing hopping term is $-\eps\sum_{i=1,2}\psi^i_b(\xi)\psi^i_b(\xi')$, with the same real functions on both copies. So the Boltzmann factor is $\cG(\xi)\cG(\xi')\exp(\sum_\iota\gamma_\iota Y_\iota(\xi)Y_\iota(\xi'))$ with all $\gamma_\iota=\eps>0$ and all $Y_\iota$ real, and for $G$ a function of one side $\theta(G)G=\overline{G(\xi')}G(\xi)$. Expanding the exponential, $Z\langle\theta(G)G\rangle$ is a sum with nonnegative coefficients of squared moduli $|\int d\xi\,G(\xi)\cG(\xi)\prod_\iota Y_\iota(\xi)^{n_\iota}|^2$, the interchange of sum and integral being justified as in \cref{lem:s07:RPbox} ($\lambda>0$, $\lmu>0$ and the growth bound). The Hermitian property $\overline{\langle\theta(G')G\rangle}=\langle\theta(G)G'\rangle$ follows from the $\theta$-invariance of the measure, so $(G,G')\mapsto\langle\theta(G')G\rangle$ is a positive semidefinite Hermitian form on bounded gauge-invariant one-sided $G$.

\emph{Step 3: $L^2$ observables.} For gauge-invariant one-sided $G$ with $\langle|G|^2\rangle<\infty$ apply Step~2 to $G_n\defeq G\mathbf 1\{|G|\le n\}$; since $|\theta(G_n)G_n|\le|\theta G||G|$ and $\langle|\theta G||G|\rangle\le\langle|G|^2\rangle$ by Cauchy--Schwarz and $\theta$-invariance, dominated convergence gives the inequality for $G=\sum_i\gamma_iG_i$, hence the form statement.

\emph{Invariance.} Under $\theta$ and under every lattice translation each term of the periodic action, in either convention, goes to a term of the same kind with the same coefficient, and the a priori measure is invariant; the value of the site-independent bare coefficient enters nowhere above.
\end{proof}

\begin{remark}[the torus paper and the literature]
(i) The torus paper \cite{SymPhase} proves \cref{prop:s10:RP} for its lattice model, King's action on the torus with a constant bare coefficient $b_T$, for every even side, every direction, both kinds of plane pair and every $b_T$ \cite[Prop.~7.9(a)]{SymPhase}; its printed scope is $\lmu\in(0,\mu_*]$, where $\mu_*=m$ is the photon mass at which that paper freezes its counterterm. The restriction plays no role in the proof, which uses no bound on $\lmu$: the two lemmas behind it there, the passage to the Proca form \cite[Lemma~7.7]{SymPhase} and reflection positivity of the Proca torus measure in a pair of planes \cite[Lemma~7.8(ii)]{SymPhase}, are stated for every photon mass $\mu>0$, with the inputs site-independent coefficients, $\lambda>0$ and even sides. So the statement lifts to every $\lmu>0$, as the proof above also shows directly. (ii) \cite[Cor.~5.4]{BFS1} (reflection positivity ``in every gauge'', for covariances $C_F=C_u+dFd^*\ge0$ with $F$ of either sign) is the same mechanism: on the torus the Feynman covariance has this form with $F=-\lmu^{-2}(\Delta_0+\lmu^2)^{-1}\le0$ on functions orthogonal to the constants. It is not a literal instance (the region of \cite[\S5]{BFS1} is a bounded subset of the lattice with a boundary condition, and reflection positivity of periodic conditions is asserted there in a remark), and \cref{prop:s10:RP} does not rest on it. (iii) The chessboard estimate uses \cref{prop:s10:RP} for functions of the $r_x$ only; nothing in this section uses reflection positivity of a cube or of a limit state.
\end{remark}

\subsection{The chessboard estimate on King's tori}

\begin{lemma}[chessboard estimate]\label{lem:s10:CB}
Let $T$ be a King torus at spacing $\eps=\eps_k$, $a$ a scheme, and $q$ a real function on the sites of $T$. In the free convention,
\begin{equation}\label{eq:s10:CB}
\ln N^a_{k,T}(0,q)\ \le\ \sum_{x}\eps^3\,\alpha^a_{k,T}(q(x)),
\end{equation}
and $u\mapsto\alpha^a_{k,T}(u)$ is convex on $\R$ with $\alpha^a_{k,T}(0)=0$. In King's torus convention $u\mapsto|T|^{-1}\ln N^{z,\mathrm{tor}}_{k,T}(0,u\mathbf 1_T)$ is likewise convex with value $0$ at $u=0$.
\end{lemma}

\begin{proof}
We apply the chessboard estimate \cite[Thm.~4.2]{FILS78} (lattice-gauge form: \cite[Thm.~2.10]{BFS1}), whose hypotheses we check: the site set of $T$ is a rectangular lattice box with sides of even length $\ell5^k$; the periodic state is invariant under translations modulo the sides and reflection positive with respect to the reflections in all planes perpendicular to the coordinate axes running midway between neighbouring sites (\cref{prop:s10:RP}; on the torus such a reflection is the reflection in a mid-plane pair, the antipodal plane $c+\ell/2$ being again a mid-plane). Take as the ``spin at $x$'' the gauge-invariant variable $r_x$ and
\[
G_x\defeq\exp\big(\eps^3q(x)(r_x^2-2\sW_\eps-a_k)\big),\qquad\text{so that}\qquad\Big\langle\prod_xG_x\Big\rangle_{k,T}=N^a_{k,T}(0,q).
\]
The $G_x$ are positive with all moments ($\lambda>0$), and the Schwarz inequality for such unbounded gauge-invariant functions, all that \cite[Thm.~4.1]{FILS78} uses beyond cyclic invariance, follows from the bounded case by truncation and monotone convergence. The conclusion is $|\langle\prod_xG_x\rangle|\le\prod_x\langle\prod_yG_x^{(y)}\rangle^{1/|\Lambda|}$, $G_x^{(y)}$ the copy of $G_x$ carried to $y$ by reflections, $|\Lambda|=|T|/\eps^3$ the number of sites. Because $\sW_\eps$, $a_k$ and $\dm_\eps$ are the same at every site, $\prod_yG_x^{(y)}=\exp(q(x)\sum_y\eps^3(r_y^2-2\sW_\eps-a_k))$ has expectation $N^a_{k,T}(0,q(x)\mathbf 1_T)=e^{|T|\alpha^a_{k,T}(q(x))}$ in the same state; taking logarithms gives \eqref{eq:s10:CB}. Convexity: $|T|\alpha^a_{k,T}(u)$ is the logarithm of a Laplace transform in the linear parameter $u$ plus a linear term, and vanishes at $u=0$; the same computation gives the convexity statement in the torus convention.
\end{proof}

\subsection{Constant sources on large tori: the extensive bound}

\begin{definition}[the extensive bound \Itwo]\label{def:s10:I2}
We say that \Itwo\ (one-sided, in the volume limit at fixed spacing; also written \UVS) holds for the scheme $a$ in the free convention if there are $s_1>0$ and $C_U\in[0,\infty)$ such that
\begin{equation}\label{eq:s10:I2}
\sup_{k\ge0}\ \limsup_{J\to\infty}\ \alpha^a_{k,T_J}(s_1)\ \le\ C_U .
\end{equation}
\end{definition}

The positive leg of the next section consumes exactly this, with $\kappa\defeq C_U/s_1$ (by convexity and $\alpha(0)=0$, $\limsup_J\alpha^a_{k,T_J}(u)\le\kappa u$ for $0\le u\le s_1$). Throughout, $M_1\in(0,m)$ is an auxiliary scalar mass, $\alpha^\infty_k(M)=M^2+\dm_k(M)$ is King's free-boundary bare coefficient on $\Lk$ at scalar mass $M$ (\cref{hyp:s03:T2}; it is the fixed-$k$ volume limit of the torus coefficient $M^2+\dm_{k,T_J}(M)$, \cref{lem:s10:V}), and, in the reference scheme $z$,
\[
s^\infty_k(M_1)\defeq\tfrac12\big(\alpha^\infty_k(m)-\alpha^\infty_k(M_1)\big),\qquad
\sigma_1=\sigma_1(M_1)\defeq\liminf_{k\to\infty}s^\infty_k(M_1),
\]
\[
B^z_{k,T}(M)\defeq E(M)-E(m)-\tfrac12\big(2c_{k,T}+z_k\big)\big(\alpha^{\mathrm{tor}}_{k,T}(m)-\alpha^{\mathrm{tor}}_{k,T}(M)\big)|T|,
\]
\[
b^\infty_k(M_1)\defeq\lim_{J\to\infty}|T_J|^{-1}B^z_{k,T_J}(M_1),\qquad
\beta_1=\beta_1(M_1)\defeq\limsup_{k\to\infty}b^\infty_k(M_1),
\]
where $\alpha^{\mathrm{tor}}_{k,T}(M)\defeq M^2+\dm_{k,T}(M)$ is the bare coefficient of King's torus convention and the limit defining $b^\infty_k$ exists by \cref{lem:s10:V}. The constant $\beta_1(M_1)$ is the second constant of \cref{hyp:s03:T2}, which says that $\sigma_1(M_1)$ and $\beta_1(M_1)$ are finite; \Tthree\ says that $\sigma_1(M_1)>0$.

\begin{proposition}[\Itwo\ from the torus inputs]\label{prop:s10:U}
Assume \Tthree\ (\cref{hyp:s03:T3}); the torus inputs are \cref{hyp:s03:T1} at the two masses $m$ and $M_1$ and \cref{hyp:s03:T2}. Then \Itwo\ holds in the reference scheme $z$ for every $s_1\in(0,\sigma_1)$, with a constant $C_U=C_U(s_1)<\infty$ depending on $(e,\lambda,m,\lmu,M)$, on $M_1$ and on $s_1$. For a scheme $a\ne z$ with $\sup_k|z_k-a_k|<\infty$ (an admissible scheme, \cref{def:s03:adm}) it holds with the constant $C_U+s_1\sup_k|z_k-a_k|$.
\end{proposition}

The proof assembles four lemmas on King's torus model at fixed spacing, stated in the form used and with their short proofs where elementary. Notation at fixed $k$ (spacing $\eps$) on $T=T_J$: $f^{\mathrm{tor}}(u)\defeq\ln N^{z,\mathrm{tor}}_{k,T}(0,u\mathbf 1)$ (King's torus model at mass $m$, Wick constant $c_{k,T}$, scheme $z$) and $f^F(u)\defeq\ln N^{z}_{k,T}(0,u\mathbf 1)$ (free convention); both are convex with $f(0)=0$ (\cref{lem:s10:CB}). Write $\alpha^{\mathrm{tor}}\defeq m^2+\dm_{k,T}(m)$, $\alpha^F\defeq m^2+\dm_\eps$, $\sigma_{k,T}\defeq\frac12(\alpha^{\mathrm{tor}}-\alpha^F)$ and $\Delta c_T\defeq\sW_\eps-c_{k,T}$ (the two \emph{convention offsets}), and $s_{k,T}(M_1)\defeq\frac12(\alpha^{\mathrm{tor}}_{k,T}(m)-\alpha^{\mathrm{tor}}_{k,T}(M_1))$.

\begin{lemma}[mass shift]\label{lem:s10:M}
If $M''>0$ satisfies $\alpha^{\mathrm{tor}}_{k,T}(M'')=\alpha^{\mathrm{tor}}_{k,T}(m)-2u$, then
\[
N^{z,\mathrm{tor}}_{k,T}(0,u\mathbf 1)=e^{B}\,Z^{(M'')}_{k,T}\big/Z^{(m)}_{k,T},\qquad B=E(M'')-E(m)-(2c_{k,T}+z_k)\,u\,|T| .
\]
\end{lemma}

\begin{proof}
The source $u\sum_x\eps^3r_x^2$ shifts the bare coefficient of the torus model at mass $m$ by $-2u$, that is to the bare coefficient of the model at mass $M''$; beyond the fixed terms, King's torus action \cite[(2.4)]{K1} contains only the counterterm and the field-independent constants $E_0,E_1$, so the two Boltzmann factors differ by the constant $e^{E(M'')-E(m)}$, and the Wick and scheme multipliers $e^{-(2c_{k,T}+z_k)u|T|}$ of $N^{z,\mathrm{tor}}$ are bookkeeping.
\end{proof}

\begin{lemma}[chord]\label{lem:s10:C}
If $f$ is convex on $\R$ with $f(0)=0$, then $f(u)\le(u/s)f(s)$ for $0\le u\le s$, $s>0$.
\end{lemma}

\begin{proof}
$u=(1-u/s)\cdot0+(u/s)\cdot s$ and $f(0)=0$.
\end{proof}

\begin{lemma}[the volume limit at fixed spacing]\label{lem:s10:V}
Fix $k$. As $J\to\infty$: $\alpha^{\mathrm{tor}}_{k,T_J}(M)\to\alpha^\infty_k(M)$ uniformly for $M\in[M_1,m]$; $s_{k,T_J}(M_1)\to s^\infty_k(M_1)$; $|T_J|^{-1}B^z_{k,T_J}(M_1)\to b^\infty_k(M_1)$, a finite limit; and the offsets vanish, $\sigma_{k,T_J}\to0$ and $\Delta c_{T_J}\to0$.
\end{lemma}

The proof is given in Appendix~\ref{app:a06_BALtorus} (\cref{ssec:a06:V}); the reason the statement holds: $\dm_{k,T}$ and $E_1$ (\cite[(3.32)--(3.33), (3.35)]{K1}) are finite sums of lattice graphs with translation-invariant massive propagators ($m,\lmu>0$), which on the torus are exact image sums of the propagators of $\eps\Z^3$, each decaying exponentially in $|x|/\eps$ at fixed $\eps$, so dominated convergence over the images gives the limits; $E_0/|T|$ is a Riemann sum of a smooth periodic log-determinant. No rate and no uniformity in the source are claimed or used. The boundedness in $J\ge J_*(k)$ of $\dm_{k,T_J}$ and $c_{k,T_J}$, used in \cref{lem:s10:F}, is a further clause of the same lemma, not one of those listed here. Of the clauses listed, Step~2 of the proof of \cref{prop:s10:U} uses the convergence of $s_{k,T_J}(M_1)$ and the $\limsup_J$ of $|T_J|^{-1}B^z_{k,T_J}(M_1)$, and Step~3 the two offset clauses.

\begin{lemma}[the fixed-spacing extensive bound]\label{lem:s10:F}
Fix $k$ and $v>0$. There is $C_k(v)<\infty$ such that for every King torus $T=T_J$ with $J\ge J_*(k)$, in either convention, $f(\pm v)\le C_k(v)|T|$; hence, by convexity, $|f(w)|\le C_k(v)\,|w|\,|T|/v$ for $|w|\le v$.
\end{lemma}

\begin{proof}
Let $I(b)\defeq\int_\C e^{-\eps^3[\frac12b|\phi|^2+\lambda|\phi|^4]}d^2\phi<\infty$ for real $b$ ($\lambda>0$), $\alpha$ the bare coefficient and $c$ the Wick constant of the convention in use, and $Z_\gamma$ the photon Gaussian integral, finite because $\lmu>0$ (at $\lmu=0$ the constant mode is in the kernel of the Feynman form). The hopping term is $\ge0$, so dropping it bounds the numerator: $\int e^{-S+w\sum\eps^3r^2}\le Z_\gamma I(\alpha-2w)^{|T|/\eps^3}$. The hopping term of a bond is $\le\eps(r_x^2+r_y^2)$, and each site lies on six bonds, so the hopping part of the action is $\le6\eps^{-2}\sum_x\eps^3r_x^2=\sum_x\eps^3\cdot\frac12(12\eps^{-2})r_x^2$ and $\int e^{-S}\ge Z_\gamma I(\alpha+12\eps^{-2})^{|T|/\eps^3}$. Hence $\ln\langle e^{w\sum\eps^3r^2}\rangle\le\eps^{-3}|T|[\ln I(\alpha-2w)-\ln I(\alpha+12\eps^{-2})]$ and
\[
f(w)\le\big\{\eps^{-3}[\ln I(\alpha-2w)-\ln I(\alpha+12\eps^{-2})]+(2|c|+|z_k|)|w|\big\}|T| .
\]
In the free convention $\alpha$ and $c$ do not depend on $T$; in the torus convention they converge as $J\to\infty$ and are bounded for $J\ge J_*(k)$ (the boundedness clause noted after \cref{lem:s10:V}). So $C_k(v)$, the supremum over the sign and over $J\ge J_*(k)$ of the brace at $w=\pm v$, is finite. The last claim is convexity with $f(0)=0$: $f(w)\le(|w|/v)f(\pm v)$ and $f(w)\ge-f(-w)\ge-(|w|/v)f(\mp v)$.
\end{proof}

\begin{proof}[Proof of \cref{prop:s10:U}]
\emph{Step 1: the exact convention bridge.} On the same torus the free-convention action differs from the torus-convention action by $\frac12(\alpha^F-\alpha^{\mathrm{tor}})\sum_x\eps^3r_x^2=-\sigma_{k,T}\sum_x\eps^3r_x^2$ plus field-independent constants. Hence $\langle e^{u\sum\eps^3r^2}\rangle^F=\langle e^{(u+\sigma)\sum\eps^3r^2}\rangle^{\mathrm{tor}}/\langle e^{\sigma\sum\eps^3r^2}\rangle^{\mathrm{tor}}$ with $\sigma=\sigma_{k,T}$, and inserting the Wick and scheme multipliers,
\begin{equation}\label{eq:s10:bridge}
f^F(u)=f^{\mathrm{tor}}(u+\sigma)-f^{\mathrm{tor}}(\sigma)-2\,\Delta c_T\,u\,|T| .
\end{equation}

\emph{Step 2: the bound in the torus convention.} Let $0\le u<\sigma_1$ and $\eta>0$. By the definitions of $\sigma_1$ and $\beta_1$ there is $K_1=K_1(u,\eta)$, which we take at least as large as the tail of \cref{hyp:s03:T1}, such that $s^\infty_k(M_1)>u$ and $b^\infty_k(M_1)<\beta_1+\eta$ for all $k\ge K_1$. Fix such a $k$. By \cref{lem:s10:V} there is $J_*\ge J_*(k)$ with $s\defeq s_{k,T_J}(M_1)>u$ for $J\ge J_*$. The chord (\cref{lem:s10:C}) at this $s$ gives $f^{\mathrm{tor}}(u)\le(u/s)f^{\mathrm{tor}}(s)\le[f^{\mathrm{tor}}(s)]_+$. \Cref{lem:s10:M} with $M''=M_1$ applies exactly, since $\alpha^{\mathrm{tor}}_{k,T_J}(M_1)=\alpha^{\mathrm{tor}}_{k,T_J}(m)-2s$ by the definition of $s$:
\[
f^{\mathrm{tor}}(s)=B^z_{k,T_J}(M_1)+\ln Z^{(M_1)}_{k,T_J}-\ln Z^{(m)}_{k,T_J}\le B^z_{k,T_J}(M_1)+\big(E_+(M_1)+E_-(m)\big)|T_J|
\]
by \cref{hyp:s03:T1}, used as an upper bound at $M_1$ and as a lower bound at $m$. Divide by $|T_J|$ and let $J\to\infty$ (\cref{lem:s10:V} for the first term):
\[
\limsup_{J\to\infty}|T_J|^{-1}f^{\mathrm{tor}}_{k,T_J}(u)\le\big[\beta_1+\eta+E_+(M_1)+E_-(m)\big]_+\qquad(k\ge K_1).
\]
For each of the finitely many $k<K_1$, \cref{lem:s10:F} with any $v\ge u$ gives $\limsup_J|T_J|^{-1}f^{\mathrm{tor}}_{k,T_J}(u)\le C_k(v)u/v<\infty$. Fix $\eta\defeq1$ (any $\eta>0$ serves); then
\[
C^{\mathrm{tor}}(u)\defeq\sup_{k\ge0}\ \limsup_{J\to\infty}|T_J|^{-1}f^{\mathrm{tor}}_{k,T_J}(u)<\infty\qquad\text{for every }u\in[0,\sigma_1).
\]

\emph{Step 3: transfer to the free convention.} Fix $u\in[0,\sigma_1)$ and $p>1$ with $pu<\sigma_1$, and let $p'=p/(p-1)$. Since $u+\sigma=\frac1p(pu)+\frac1{p'}(p'\sigma)$, convexity gives $f^{\mathrm{tor}}(u+\sigma)\le\frac1pf^{\mathrm{tor}}(pu)+\frac1{p'}f^{\mathrm{tor}}(p'\sigma)$, whatever the sign of $\sigma=\sigma_{k,T_J}$. Once $p'|\sigma_{k,T_J}|\le v$, which holds for all large $J$ by \cref{lem:s10:V}, \cref{lem:s10:F} gives $|f^{\mathrm{tor}}(p'\sigma)|,|f^{\mathrm{tor}}(\sigma)|\le C_k(v)p'|\sigma_{k,T_J}||T_J|/v$. Inserting this into \eqref{eq:s10:bridge},
\[
|T_J|^{-1}f^F_{k,T_J}(u)\ \le\ \tfrac1p|T_J|^{-1}f^{\mathrm{tor}}_{k,T_J}(pu)+\tfrac{2C_k(v)p'}{v}|\sigma_{k,T_J}|+2|\Delta c_{T_J}|\,u
\ \xrightarrow[J\to\infty]{\ \limsup\ }\ \tfrac1pC^{\mathrm{tor}}(pu),
\]
because $\sigma_{k,T_J}\to0$ and $\Delta c_{T_J}\to0$ at fixed $k$ (\cref{lem:s10:V}; the factor $\frac1{p'}\le1$ in the middle term has been dropped). Hence $\sup_k\limsup_J\alpha^z_{k,T_J}(u)\le\frac1pC^{\mathrm{tor}}(pu)<\infty$ for every $u\in[0,\sigma_1)$. Take $u=s_1\in(0,\sigma_1)$ and $C_U\defeq[\inf_{p>1:\,ps_1<\sigma_1}\frac1pC^{\mathrm{tor}}(ps_1)]_+$; this is \eqref{eq:s10:I2} in the scheme $z$. For a scheme $a\ne z$, $N^a_{k,T}(0,q)=e^{-(a_k-z_k)\sS(q)}N^z_{k,T}(0,q)$ gives $\alpha^a_{k,T}(u)=\alpha^z_{k,T}(u)+(z_k-a_k)u\le\alpha^z_{k,T}(u)+|z_k-a_k|\,u$, so \eqref{eq:s10:I2} holds with $C_U+s_1\sup_k|z_k-a_k|$.
\end{proof}

\begin{proposition}[the auxiliary mass on either side of $m$]\label{prop:s10:Upm}
Let $M_1>0$, $M_1\ne m$. Assume \cref{hyp:s03:T1} at $m$ and at $M_1$, and \cref{hyp:s03:T2} and \Tthree\ read at this $M_1$: $s^\infty_k(M_1)=\frac12(\alpha^\infty_k(m)-\alpha^\infty_k(M_1))$, of either sign, has $\sigma_1(M_1)=\liminf_ks^\infty_k(M_1)>0$, and $\beta_1(M_1)<\infty$. Then \Itwo\ holds in the scheme $z$ for every $s_1\in(0,\sigma_1(M_1))$, and for $a\ne z$ with the shifted constant, exactly as in \cref{prop:s10:U}.
\end{proposition}

\begin{proof}
The proof of \cref{prop:s10:U} is repeated with \cref{lem:s10:V} stated on the interval $I\defeq[\min(M_1,m),\max(M_1,m)]$, and we check that no step uses $M_1<m$. \Cref{lem:s10:M} is an identity for any $M''$ with $\alpha^{\mathrm{tor}}(M'')=\alpha^{\mathrm{tor}}(m)-2u$, whatever the order of $M''$ and $m$. \Cref{lem:s10:C} needs $0\le u\le s$ only. \Cref{lem:s10:V}: its proof uses only that the propagators are translation-invariant massive lattice covariances at fixed $\eps$ with masses in $I\cup\{\lmu\}$, bounded below by $\min(M_1,m,\lmu)>0$, so it holds on $I$ word for word; Step~2 uses it at the two masses $m$ and $M_1$ only. \Cref{lem:s10:F} carries no mass order, and \eqref{eq:s10:bridge} is an identity. In Step~2 the sign clause is $s^\infty_k(M_1)>u\ge0$, the chord is applied at $s=s_{k,T_J}(M_1)>u\ge0$, \cref{lem:s10:M} at $M''=M_1$ is exact by the definition of $s$, and \cref{hyp:s03:T1} is used as an upper bound at $M_1$ and a lower bound at $m$; none of this uses $M_1<m$. Step~3 carries no mass.
\end{proof}

The two-sided variant matters exactly on the set $E^<_W\setminus E_W$ of couplings at which $m$ minimises $\alpha_\infty$ on $(0,m]$ but not on $(0,\infty)$ (\cref{lem:s14:CW}; for instance when $\alpha'_\infty(m)<0$ and $m$ minimises $\alpha_\infty$ on $(0,m]$): there \Tthree\ fails for every $M_1<m$ while some $M_1>m$ has $\sigma_1(M_1)>0$. The main results use \cref{prop:s10:U} with $M_1\in(0,m)$, as \Tthree\ is posed.

\begin{remark}[the convention of \cref{hyp:s03:T1} and how much of it is used]\label{rem:s10:T1convention}
(i) \Cref{hyp:s03:T1} is posed in King's torus convention (counterterm and Wick constant computed with the periodic propagators of the torus itself), not in the normalisation of Bałaban's own papers. This is what the proof needs: the offsets $\sigma_{k,T_J}$, $\Delta c_{T_J}$ between King's torus convention and the free convention vanish in the volume limit at fixed spacing (\cref{lem:s10:V}), and Step~3 of the proof of \cref{prop:s10:U} transfers the bound through them; between other normalisations they need not vanish. (ii) Only a tail of \cref{hyp:s03:T1} is used: Step~2 of that proof invokes it for $k\ge K_1$ and $J\ge J_*(k)$, the finitely many $k<K_1$ being handled by \cref{lem:s10:F} at fixed spacing, and the constants $E_\pm(M)$ are independent of $k$ and of the torus. (iii) In the scheme $z$, $C_U$ depends on the couplings, on $M$ and on the auxiliary choices $M_1$, $s_1$, and on nothing else; neither it nor any other bound of this section is uniform in $\lmu$.
\end{remark}
